\justifying
\NSFCBodyText

\subsubsection{年度研究计划}

\textbf{第一年：可靠证据状态与运输风险辨识。} 完成事件、三类时间戳、模态质量、缓存/重传和标签的数据规范；同步启动运输与仓储数据准备，建立可控异常、弱网回放、强基线和泄漏审计。完成三时钟观测模型、敏感性区间与拒判方法，形成事件风险、证据完整度和不确定性输出；以预实验和功效分析冻结主要终点与统计计划。年度门控为可观测锚点、三时钟日志和关键分层支持达到预设标准；未通过时只报告部分辨识与拒判，不主张点风险可辨识。

\textbf{第二年：集合条件价值与弱网安全协同。} 研究证据互补、冗余、交付概率和迟到失效下的价值估计及置信下界；完成采样/压缩—缓存/传输两级决策、边缘计算和安全屏障。在相同 2G/3G 模拟轨迹与资源预算下完成价值校准、成对重放、断点续传、可靠上传和低功耗评价；同步建立仓储本地风险与漂移基线。年度门控为母流覆盖、动作支持和资源记录可追溯；未通过时限于受控重放，不从单路径日志外推未执行动作价值。

\textbf{第三年：仓储校准预警与有限闭环。} 完成温湿度、烟雾、门禁和视频的多时间尺度风险校准、异常识别、分级预警、拒判和应急响应建议。仓储本地任务独立验收；统一货物标识、实体配对和统计功效同时达标后，才检验事件特异运输记忆。完成运输感知—弱网保全—仓储预警的离线/半实物闭环和误差传播分析，明确适用范围与失败条件。

\subsubsection{预期研究结果}

形成三时钟弱网多模态风险辨识方法、事件证据集合价值与安全资源决策方法、漂移环境下仓储风险校准及条件性预测记忆方法；交付可复现的数据与三时钟日志规范、弱网回放及评测代码、边缘原型和反证实验报告，并据此形成论文、软件或数据成果。

\subsubsection{学术交流与合作计划}

围绕多模态时序异常、边缘智能和智能物流，参加相关国内外学术会议并开展阶段成果交流；结合数据、设备和场景条件的落实情况，与校内外相关团队就网络回放、端侧测量和物流场景验证开展合作。
